\documentclass[10pt,a4paper]{amsart}
\usepackage[margin=19mm]{geometry}
\usepackage{amsmath,amssymb,mathtools}
\usepackage[scr=rsfs,cal=euler]{mathalfa}
\usepackage{xfrac}
\usepackage[unicode,hidelinks,bookmarks=false]{hyperref}
\newcommand{\ie}{i.e.\ }
\newcommand{\cf}{cf.\ }
\newcommand{\resp}{respectively\ }
\newcommand{\Subsection}{\subsection}
\providecommand{\nobreakdash}{\nobreak\hbox{-}}
\setlength{\emergencystretch}{2em}
\multlinegap0pt
\usepackage{bm}
\usepackage{bbm}
\usepackage{dsfont}
\usepackage{xspace}
\usepackage{cancel}
\usepackage{mathtools}
\usepackage{amsthm}
\makeatletter
\@ifundefined{proposition}{\newtheorem{proposition}{Proposition}}{}
\@ifundefined{theorem}{\newtheorem{theorem}{Theorem}}{}
\makeatother
\title[A Neurofinance Framework for Subjective Temporal Perception,]{A Neurofinance Framework for Subjective Temporal Perception, Risk, and Investment Behavior}
\author{Pascal Stiefenhofer}
\date{Source: arXiv:2608.14930v1 (CC BY 4.0)}
\begin{document}
\maketitle

\noindent\emph{Source and licence.} A Neurofinance Framework for Subjective Temporal Perception, Risk, and Investment Behavior. Version arXiv:2608.14930v1 (CC BY 4.0), distributed under the Creative Commons Attribution 4.0 International licence, as recorded in the arXiv licence metadata for this exact version and shown on the abstract page https://arxiv.org/abs/2608.14930v1. Full rights notice: licenses/arxiv-2608.14930-CC-BY-4.0.txt.\par\smallskip

\noindent\textit{Editorial scope.} Counted unit: the Theorem whose source label is \texttt{thm:dispersion} in the article (source0.tex lines 1095--1151, source label \texttt{thm:dispersion}), delivered together with the article's own complete original proof of it (source0.tex lines 1153--1262, one \texttt{proof} environment of 110 lines). No mathematics is written, repaired or re-derived: statement and proof are the article's own text line for line. Required context retained uncounted: prop:temporal-process (lines 736--765). Every definition, display and label of those lines is retained unchanged. Omitted from this package and not redistributed: 1--735, 766--1094, 1152--1152, 1263--2438. Nothing inside a retained excerpt is omitted or altered: every retained body is delivered exactly as the article's lines are written, so no omission receipt of any kind accompanies this package. Citations: the retained excerpts carry no citation command at all, so no bibliography entry of the article is transcribed and no reference is redistributed. Reference adaptation: none was needed and none was made; every reference inside the delivered unit resolves inside it, so no unresolved reference or citation is produced. Printed slips kept literal: no printed word, symbol, hypothesis or number of the counted statement or of its proof was altered. Layout and class: the article's own class is \texttt{SNmult} with packages including type1cm, makeidx, graphicx, multicol, footmisc, newtxtext, newtxmath, amsmath, mathtools, bm, bbm, mathrsfs, booktabs, longtable; the wrapper is the lane's pinned \texttt{amsart} editorial wrapper at 10pt on A4, which restates the theorem-like environments the retained text needs and adds the article's own macro definitions that the retained text needs (none), with extra wrapper packages (bm, bbm, dsfont, xspace, cancel, mathtools). The delivered numbering is the unit's own. Assets: no retained span contains a figure environment, a graphics-inclusion command, a PDF-page inclusion, a table or a TikZ picture; no image file is copied into this package, the package declares no asset dependency and no image file is redistributed with it. Licence: version arXiv:2608.14930v1 of this article is distributed under the Creative Commons Attribution 4.0 International licence, as recorded in the arXiv licence metadata for this exact version and shown on the abstract page https://arxiv.org/abs/2608.14930v1. A full rights notice is delivered at licenses/arxiv-2608.14930-CC-BY-4.0.txt. Characters: every retained excerpt is pure ASCII, so the delivered engine, XeLaTeX with its default Latin Modern text font, reports no missing character.
\begin{proposition}[Temporal process induced by admissible evolution]
\label{prop:temporal-process}

Every admissible history
\(X=(X_F,X_E)\in\mathcal S_T(x_0)\), together with an initial temporal
normalisation \(\tau_0\in\mathbb R\), determines a unique absolutely continuous
subjective financial-time process
\[
\tau_X(t)
=
\tau_0+
\int_0^t
\omega_{X(s)}\!\left(\dot X(s)\right)\,ds.
\]
Moreover, \(\tau_X\) is non-decreasing and
\begin{equation}
\dot\tau_X(t)
=
\omega_{X(t)}\!\left(\dot X(t)\right)
=
\omega_{F,X(t)}\!\left(\dot X_F(t)\right)
+
\omega_{E,X(t)}\!\left(\dot X_E(t)\right)
\in
\Lambda(X(t))
\quad
\text{for a.e. }t\in[0,T].
\label{eq:temporal-process-rate}
\end{equation}
\end{proposition}

\begin{theorem}[Dynamic representation of the admissible rate spectrum]
\label{thm:dispersion}

Let \(x_0\in M\), and suppose that \(A\) is locally
\(C^1\)-velocity-realisable at \(x_0\). Then
\begin{equation}
\Lambda(x_0)
=
\left\{
\dot\tau_X(0):
X\in\mathcal R^1(x_0)
\right\}.
\label{eq:representation-spectrum}
\end{equation}
Consequently,
\begin{equation}
\Delta_\tau(x_0)
=
\operatorname{diam}
\left\{
\dot\tau_X(0):
X\in\mathcal R^1(x_0)
\right\}.
\label{eq:representation-dispersion}
\end{equation}

Moreover, there exist
\(X_-,X_+\in\mathcal R^1(x_0)\) such that
\[
\dot\tau_{X_-}(0)
=
\underline\lambda(x_0),
\qquad
\dot\tau_{X_+}(0)
=
\overline\lambda(x_0),
\]
and, under the common temporal normalisation
\(\tau_{X_-}(0)=\tau_{X_+}(0)=\tau_0\),
\begin{equation}
\tau_{X_+}(t)-\tau_{X_-}(t)
=
\Delta_\tau(x_0)t+o(t)
\qquad
(t\downarrow0).
\label{eq:dispersion-expansion}
\end{equation}
Hence
\[
\Delta_\tau(x_0)
=
\lim_{t\downarrow0}
\frac{
\tau_{X_+}(t)-\tau_{X_-}(t)
}{t}.
\]
\end{theorem}

\begin{proof}

Let \(X\in\mathcal R^1(x_0)\). Since
\(\dot X(0)\in A(x_0)\), Proposition~\ref{prop:temporal-process} gives
\[
\dot\tau_X(0)
=
\omega_{x_0}\!\left(\dot X(0)\right)
\in\Lambda(x_0).
\]
Hence
\[
\left\{
\dot\tau_X(0):
X\in\mathcal R^1(x_0)
\right\}
\subseteq
\Lambda(x_0).
\]
Conversely, let \(\lambda\in\Lambda(x_0)\). By
\(\Lambda(x_0)=\omega_{x_0}(A(x_0))\), there exists
\(v\in A(x_0)\) such that
\[
\lambda=\omega_{x_0}(v).
\]
Local \(C^1\)-velocity realisability at \(x_0\) yields
\(X\in\mathcal R^1(x_0)\) with \(\dot X(0)=v\). Therefore,
Proposition~\ref{prop:temporal-process} gives
\[
\dot\tau_X(0)
=
\omega_{x_0}(v)
=
\lambda.
\]
Thus
\[
\Lambda(x_0)
\subseteq
\left\{
\dot\tau_X(0):
X\in\mathcal R^1(x_0)
\right\},
\]
which proves \eqref{eq:representation-spectrum}. Taking diameters gives
\eqref{eq:representation-dispersion}. Because \(A(x_0)\) is compact and
\(v\mapsto\omega_{x_0}(v)\) is continuous, there exist
\(v_-,v_+\in A(x_0)\) satisfying
\[
\omega_{x_0}(v_-)
=
\underline\lambda(x_0),
\qquad
\omega_{x_0}(v_+)
=
\overline\lambda(x_0).
\]
Local \(C^1\)-velocity realisability provides
\(X_-,X_+\in\mathcal R^1(x_0)\) with
\(\dot X_-(0)=v_-\) and \(\dot X_+(0)=v_+\). Hence
\[
\dot\tau_{X_-}(0)
=
\underline\lambda(x_0),
\qquad
\dot\tau_{X_+}(0)
=
\overline\lambda(x_0).
\]
Define
\[
r_\pm(t)
=
\omega_{X_\pm(t)}
\!\left(\dot X_\pm(t)\right).
\]
Since \(X_\pm\) are \(C^1\) and \(\omega\) is continuous,
\(r_\pm\) are continuous at \(t=0\). Therefore
\[
r_\pm(t)=r_\pm(0)+o(1)
\qquad
(t\downarrow0),
\]
and hence
\[
r_+(t)-r_-(t)
=
\Delta_\tau(x_0)+o(1).
\]
Using the common temporal normalisation,
\[
\begin{aligned}
\tau_{X_+}(t)-\tau_{X_-}(t)
&=
\int_0^t
\left[r_+(s)-r_-(s)\right]\,ds \\
&=
\Delta_\tau(x_0)t+o(t),
\end{aligned}
\]
which proves \eqref{eq:dispersion-expansion}. Dividing by \(t>0\) and taking
\(t\downarrow0\) gives
\[
\Delta_\tau(x_0)
=
\lim_{t\downarrow0}
\frac{\tau_{X_+}(t)-\tau_{X_-}(t)}{t},
\]
completing the proof.
\end{proof}

\end{document}
